% Part: normal-modal-logic
% Chapter: filtrations
% Section: preliminaries

\documentclass[../../../include/open-logic-section]{subfiles}

\begin{document}

\olfileid[id]{nml}{fil}{pre}

\olsection{Persiapan}

Filtrasi memungkinkan kita menetapkan bahwa sistem-sistem logika modal kita
dapat diputuskan dengan menunjukkan bahwa sistem tersebut mempunyai
\emph{sifat model berhingga}, yakni bahwa setiap !!{formula} yang bernilai
benar (salah) dalam suatu model juga bernilai benar (salah) dalam suatu model
\emph{berhingga}. Filtrasi didefinisikan relatif terhadap himpunan
!!{formula}s yang tertutup terhadap subformula.

\begin{defn}\ollabel{defn:modallyclosed}
  Suatu himpunan !!{formula}s~$\Gamma$ \emph{tertutup terhadap subformula}
  jika himpunan itu memuat setiap subformula dari setiap !!a{formula} dalam
  $\Gamma$. Selanjutnya, $\Gamma$ \emph{tertutup secara modal} jika himpunan
  itu tertutup terhadap subformula dan, selain itu, $!A \in \Gamma$
  mengimplikasikan $\Box!A, \Diamond!A \in \Gamma$.
\end{defn}

Sebagai contoh, jika diberikan !!a{formula}~$!A$, himpunan semua
sub-!!{formula}snya tertutup terhadap sub-!!{formula}s. Ketika kita
mendefinisikan filtrasi suatu model melalui himpunan sub-!!{formula}s
dari~$!A$, filtrasi tersebut akan mempunyai sifat yang kita cari: filtrasi itu
membuat $!A$ bernilai benar (salah) jika dan hanya jika model semula demikian.

Himpunan dunia suatu filtrasi dari~$\mModel{M}$ melalui~$\Gamma$ didefinisikan
sebagai himpunan semua kelas ekuivalensi dari relasi ekuivalensi berikut.

\begin{defn}
Misalkan $\mModel{M} =\tuple{W, R, V}$ dan andaikan $\Gamma$ tertutup terhadap
sub-!!{formula}s. Definisikan suatu relasi $\equiv$ pada $W$ yang berlaku bagi
sebarang dua dunia yang membuat !!{formula}s yang sama dari $\Gamma$ bernilai
benar, yakni:
\[
u \equiv v \quad \text{jika dan hanya jika} \quad
\forall !A \in \Gamma : \mSat{M}{!A}[u] \Leftrightarrow \mSat{M}{!A}[v].
\]
Kelas ekuivalensi~$[w]_\equiv$ dari suatu dunia~$w$, atau disingkat $[w]$,
adalah himpunan semua dunia yang $\equiv$-ekuivalen dengan~$w$:
\[
[w] = \Setabs{v}{v \equiv w}.
\]
\end{defn}

\begin{prop}
  Jika diberikan $\mModel{M}$ dan $\Gamma$, relasi $\equiv$ yang didefinisikan
  di atas merupakan relasi ekuivalensi, yakni bersifat refleksif, simetris,
  dan transitif.
\end{prop}

\begin{proof}
  Relasi $\equiv$ bersifat refleksif, sebab $w$ membuat tepat
  !!{formula}s yang sama dari~$\Gamma$ bernilai benar seperti dirinya sendiri.
  Relasi itu bersifat simetris, sebab jika $u$ membuat !!{formula}s yang sama
  dari~$\Gamma$ bernilai benar seperti $v$, hal yang sama berlaku bagi $v$
  dan~$u$. Relasi itu juga bersifat transitif, sebab jika $u$ membuat
  !!{formula}s yang sama dari~$\Gamma$ bernilai benar seperti~$v$, dan $v$
  seperti $w$, maka $u$ membuat !!{formula}s yang sama dari~$\Gamma$ bernilai
  benar seperti~$w$.
\end{proof}

Relasi $\equiv$, sebagaimana setiap relasi ekuivalensi, mempartisi~$W$ menjadi
\emph{kelas-kelas ekuivalensi}, yakni himpunan bagian dari~$W$ yang saling lepas
berpasangan dan bersama-sama mencakup seluruh~$W$. Setiap $w \in W$ merupakan
!!a{element} dari salah satu kelas tersebut, yaitu~$[w]$, sebab $w \equiv w$.
Jadi, kelas-kelas $[w]$ mencakup seluruh~$W$. Kelas-kelas itu saling lepas
berpasangan, sebab jika $u \in [w]$ dan $u \in [v]$, maka $u \equiv w$ dan
$u \equiv v$; berdasarkan simetri dan transitivitas, $w \equiv v$, sehingga
$[w] = [v]$.

\end{document}
